\documentclass[11pt]{article}
\usepackage[a4paper,margin=25mm]{geometry}
\usepackage{amsmath,amssymb,amsthm,booktabs,microtype}
\usepackage[hidelinks]{hyperref}

\newtheorem{theorem}{Theorem}
\newtheorem{lemma}{Lemma}

\title{A Symmetric Six-Point Counterexample to Concavity of Determinantal Entropy}
\author{RenZhenzhuo\\\small on behalf of the Euler Agent Team}
\date{Review note, 6 September 2026}

\begin{document}
\maketitle

\begin{abstract}
Lyons asked whether the Shannon entropy of a finite determinantal probability
measure is concave in its positive-contraction kernel.  We give a counterexample
on six labelled points.  Its midpoint is a real symmetric matrix obtained by
moving a rank-three projection into the interior, and its two endpoints differ
from the midpoint by opposite purely imaginary Hermitian matrices.  Symmetry
reduces the 64 complete-event probabilities to eight pairs.  A short exact
calculation proves that the endpoint laws agree and that their entropy is
strictly larger than the entropy at the midpoint.
\end{abstract}

\section{The example}

Let $E=[6]$.  For a Hermitian matrix $L$ with $0\preceq L\preceq I$,
write $X_L$ for the determinantal process characterized by
\[
 \Pr(T\subseteq X_L)=\det L[T]\qquad(T\subseteq E),
\]
and write
\[
 p_A(L)=\Pr(X_L=A),\qquad
 H(L)=-\sum_{A\subseteq E}p_A(L)\log p_A(L).
\]
The conjecture that $H$ is concave in $L$ was posed by
Lyons~\cite{Lyons2003DPM} and restated as Conjecture~2.6 in his ICM
survey~\cite{Lyons2014}.

\begin{theorem}\label{thm:main}
There are strict positive contractions $Q_+$ and $Q_-$ whose midpoint $K$
satisfies
\[
 p_A(Q_+)=p_A(Q_-)\qquad(A\subseteq E)
\]
and
\begin{equation}\label{eq:gap}
 \frac1{40\,000\,000}
 <H(Q_+)-H(K)
 <\frac1{39\,000\,000}.
\end{equation}
Consequently,
\[
 H(K)<\frac{H(Q_-)+H(Q_+)}2.
\]
Thus determinantal entropy is not concave on the set of complex Hermitian
positive contractions.
\end{theorem}

We now give the matrices.  Consider
the two $3\times3$ matrices
\[
 R=\frac13
 \begin{pmatrix}
  2&-1& 2\\
 -1& 2& 2\\
  2& 2&-1
 \end{pmatrix},
 \qquad
 S=\begin{pmatrix}
 0&1&-1\\
 -1&0&1\\
 1&-1&0
 \end{pmatrix},
\]
and put
\begin{equation}\label{eq:PB}
 P=\frac12\begin{pmatrix}I_3&R\\ R&I_3\end{pmatrix},
 \qquad
 B=\begin{pmatrix}S&0\\0&-S\end{pmatrix}.
\end{equation}
Finally, set
\begin{equation}\label{eq:kernels}
 \varepsilon=\frac1{475},\qquad
 t=\frac1{2300},\qquad
 K=\varepsilon I_6+(1-2\varepsilon)P,
 \qquad Q_\pm=K\pm itB.
\end{equation}

Numerically, the entropy increase in \eqref{eq:gap} is
$2.517182345129681\ldots\times10^{-8}$.  For scale,
\[
 H(K)=2.877726131217012\ldots,\qquad
 H(Q_+)=2.877726156388836\ldots;
\]
the increase is about $8.75\times10^{-9}$ times the entropy at the midpoint.

\section{The structural checks}

Everything needed to see that the matrices in \eqref{eq:kernels} are legal is
contained in the following identities.

\begin{lemma}\label{lem:structure}
The matrices in \eqref{eq:PB} satisfy
\[
 R^T=R,\qquad R^2=I_3,\qquad
 S^T=-S,\qquad RS=-SR.
\]
Hence $P$ is a rank-three orthogonal projection, $B^T=-B$, and $PB=BP$.
Furthermore, $\lVert iB\rVert=\sqrt3$.
\end{lemma}

\begin{proof}
The four displayed identities follow by direct multiplication.  Block
multiplication then gives $P^T=P=P^2$, $B^T=-B$, and $PB=BP$.  The
characteristic polynomial of $iS$ is
$\lambda(\lambda^2-3)$, which gives the norm assertion.
\end{proof}

The eigenvalues of $K$ are $\varepsilon$ and $1-\varepsilon$, each with
multiplicity three.  Since
\[
 t\lVert iB\rVert=t\sqrt3<\varepsilon,
\]
the usual eigenvalue perturbation bound places every eigenvalue of $Q_\pm$
strictly between $0$ and $1$.  Thus $Q_\pm$ are strict positive
contractions.

For later use, recall the complete-event formula
\begin{equation}\label{eq:event}
 p_A(L)=(-1)^{|E\setminus A|}
 \det\bigl(L-I_{E\setminus A}\bigr),
\end{equation}
where $I_C$ denotes the diagonal projection onto the coordinates in $C$.
It follows immediately from inclusion--exclusion applied to the defining
inclusion probabilities.

Because $Q_-=\overline{Q_+}$, the two determinants in \eqref{eq:event} are
complex conjugates.  They are real because their matrices are Hermitian.
Therefore
\begin{equation}\label{eq:same-law}
 p_A(Q_+)=p_A(Q_-)\qquad(A\subseteq E).
\end{equation}
This proves the equality of the endpoint laws without numerical computation.

\section{The entropy comparison}

It remains to establish the strict inequality in \eqref{eq:gap}.  The following
table contains the entire calculation.  Put
\[
 u=\varepsilon(1-\varepsilon).
\]
Coordinate permutations preserving $P$ and sending $B$ to either $B$ or
$-B$, together with complementation, divide the 64 subsets into the eight
classes shown below.  For example, the coordinate symmetries are generated by
$(1\,2\,3)(4\,6\,5)$, $(2\,3)(4\,6)$, and
$(1\,4)(2\,5)(3\,6)$.  Conjugation by
$\mathrm{diag}(I_3,-I_3)$ sends $Q_-$ to $I-Q_+$; the standard
complementation rule for DPPs and \eqref{eq:same-law} therefore identify a
subset with its complement.  Substitution into \eqref{eq:event} gives the two
right-hand columns.

Here is one row in full.  Since $K$ has eigenvalues $\varepsilon$ and
$1-\varepsilon$, each three times,
\[
 p_{\varnothing}(K)=\det(I-K)
 =\varepsilon^3(1-\varepsilon)^3=u^3.
\]
The commutation $PB=BP$ makes both the range and the kernel of $P$ invariant
under $iB$.  Under the identifications
$x\mapsto(x,Rx)$ and $x\mapsto(x,-Rx)$, the restriction of $iB$ to either
three-dimensional subspace is represented by $iS$, whose eigenvalues are
$0,\sqrt3,-\sqrt3$.  Therefore
\begin{align*}
 p_{\varnothing}(Q_+)
 &=\det(I-Q_+)\\
 &=\varepsilon(1-\varepsilon)
   (\varepsilon^2-3t^2)((1-\varepsilon)^2-3t^2)\\
 &=u^3+3t^2u(3t^2+2u-1).
\end{align*}
This is exactly the first row of the table.  The other seven rows follow from
the same determinant formula; the checker verifies every representative and
every multiplicity.

\begin{center}
\setlength{\tabcolsep}{3.5pt}
\renewcommand{\arraystretch}{1.25}
\scriptsize
\begin{tabular}{@{}ccll@{}}
\toprule
representative $A$&multiplicity&$p_A(K)$&$p_A(Q_+)-p_A(K)$\\
\midrule
\(\varnothing\)&2
 &\(u^3\)
 &\(3t^2u(3t^2+2u-1)\)\\
\(\{1\}\)&12
 &\(u^2(1-2u)/2\)
 &\(-t^2(18t^2u-3t^2+12u^2-6u+1)/2\)\\
\(\{1,2\}\)&12
 &\(u(1-2u)^2/4\)
 &\(t^2(18t^2u-6t^2+12u^2-8u+1)/2\)\\
\(\{1,2,3\}\)&2
 &\((1-2u)^3/8\)
 &\(3t^2(1-2u)(3t^2+2u-1)/2\)\\
\(\{1,4\}\)&12
 &\(u(36u^2-20u+5)/36\)
 &\(t^2(54t^2u-12t^2+36u^2-10u+1)/6\)\\
\(\{2,4\}\)&6
 &\(u(9u^2-8u+2)/9\)
 &\(t^2(27t^2u-6t^2+18u^2-11u+2)/3\)\\
\(\{1,2,4\}\)&12
 &\((1-2u)(9u^2-4u+1)/18\)
 &\(-t^2(54t^2u-15t^2+36u^2-20u+2)/6\)\\
\(\{1,3,4\}\)&6
 &\((1-2u)(36u^2-4u+1)/72\)
 &\(-t^2(54t^2u-15t^2+36u^2-8u-1)/6\)\\
\bottomrule
\end{tabular}
\end{center}

For completeness, we record how the numerical sign is certified without
floating-point assumptions.  If a positive rational number is written as
$x=2^kr$, where $1\le r<2$, and $y=(r-1)/(r+1)$, then for every positive
integer $m$,
\begin{equation}\label{eq:log-bound}
 2\sum_{j=0}^{m-1}\frac{y^{2j+1}}{2j+1}
 \le \log r\le
 2\sum_{j=0}^{m-1}\frac{y^{2j+1}}{2j+1}
 +\frac{2y^{2m+1}}{(2m+1)(1-y^2)}.
\end{equation}
The same formula at $r=2$ gives rational lower and upper bounds for
$\log2$.  Thus \eqref{eq:log-bound} gives directed rational bounds for the
logarithm of every entry in the table.

For a row of multiplicity $m$, with center probability $p$ and endpoint change
$d$, its contribution to the entropy difference is
$m[-(p+d)\log(p+d)+p\log p]$.  Substituting
$\varepsilon=1/475$ and $t=1/2300$, summing the eight rows, and taking 50
terms in \eqref{eq:log-bound} gives, by rational arithmetic,
\[
 \frac1{40\,000\,000}
 <H(Q_+)-H(K)
 <\frac1{39\,000\,000}.
\]
This is \eqref{eq:gap}.  Together with \eqref{eq:same-law}, it completes the
proof of Theorem~\ref{thm:main}.

The accompanying checker starts from the four displayed matrices and
reproduces every event probability, the eight-row grouping,
endpoint feasibility, equality of the endpoint laws, and the two directed
entropy bounds using only exact rational arithmetic and the Python standard
library.

\section{Why the direction works}

The calculation above also has a simple conceptual explanation.  If $L$ is
real symmetric and $B$ is real skew-symmetric, transposition in
\eqref{eq:event} shows that
\[
 p_A(L+isB)=p_A(L-isB).
\]
Every complete-event probability is therefore an even function of $s$, so
its first derivative at $s=0$ vanishes.  Twice differentiating entropy gives
\begin{equation}\label{eq:second-derivative}
 \left.\frac{d^2}{ds^2}H(L+isB)\right|_{s=0}
 =-\sum_{A\subseteq E}
 \left.\frac{d^2}{ds^2}p_A(L+isB)\right|_{s=0}\log p_A(L).
\end{equation}
The normally negative term involving the squares of the first derivatives has
disappeared.  Formula \eqref{eq:second-derivative} says that the sign is decided
entirely by the second-order motion of the event probabilities.  In
information-geometric language, this is the second fundamental form of the
probability map paired with the entropy gradient.

The two displayed matrices are one member of a simple algebraic family.  If
$R$ is any real symmetric matrix with $R^2=I$, and $S$ is any real
skew-symmetric matrix satisfying $RS=-SR$, then the block formulas
\[
 P=\frac12\begin{pmatrix}I&R\\R&I\end{pmatrix},
 \qquad
 B=\begin{pmatrix}S&0\\0&-S\end{pmatrix}
\]
automatically give an orthogonal projection $P$ and a commuting skew direction
$B$.  The particular rational pair above was selected within this family
because its event probabilities have only eight classes and its second-order
motion raises entropy.

This family was found by searching for first-order degeneracy near the boundary
of the set of positive contractions.  That search motivates the construction
but is not needed to verify it.  Once $R$ and $S$ are given, the proof consists
of the matrix identities in Lemma~\ref{lem:structure}, the event formula
\eqref{eq:event}, and the eight-row entropy calculation.

The midpoint in the example is real symmetric, whereas the direction $iB$ is
purely imaginary Hermitian.  The example therefore disproves the conjecture
with complex Hermitian kernels.  It does not settle the separate restriction
to real symmetric kernels, nor does it establish that six is the smallest
possible dimension.

The construction was found with Qwen Euler agent assistance.  Numerical search
was used to locate the pattern; all statements in this note and in the checker
are verified by exact arithmetic once the matrices $R$ and $S$ are given.
To the best of our knowledge, the cited conjecture had no previously published
proof or counterexample.

\bibliographystyle{plain}
\bibliography{references}

\end{document}
